\documentclass[11pt,a4paper]{article}
\usepackage[utf8]{inputenc}
\usepackage[T1]{fontenc}
\usepackage[french]{babel}
\usepackage[left=2cm,right=2cm,top=2cm,bottom=2cm]{geometry}
\usepackage{hyperref}
\usepackage{fontawesome5}

\pagestyle{empty}
\setlength{\parindent}{0pt}
\setlength{\parskip}{1em}

\begin{document}

% --- En-tête ---
\begin{center}
    {\Large \textbf{Loïc Aron MBASSI EWOLO}} \\
    \vspace{0.1cm}
    \small
    \faEnvelope\ \href{mailto:loicaron.mbassiewolo@ensea.fr}{loicaron.mbassiewolo@ensea.fr} \quad
    \faPhone\ +33 7 59 52 33 98 \quad
    \faGithub\ \href{https://github.com/Nameless0l}{Nameless0l}
\end{center}

\vspace{0.5cm}

\textbf{Objet :} Candidature Stage - Contribution à un environnement d'aide à l'ingénierie basé sur l'IA

Madame, Monsieur,

Je candidate pour rejoindre Airbus Connected Intelligence sur le sujet de l'aide à l'ingénierie système par l'IA. Ce stage m'intéresse car il cible précisément mon domaine d'expertise : concevoir des outils intelligents qui assistent les ingénieurs dans leurs processus métier.

Mon profil répond concrètement à vos missions pour trois raisons :

\textbf{1. Je sais construire des générateurs de code et outils d'aide à l'ingénierie.}
Votre offre demande de définir et mettre en œuvre des environnements IA. J'ai déjà réalisé ce travail avec mon projet \textbf{Laravel API Generator} (open source) : un outil qui génère automatiquement une architecture backend complète depuis des diagrammes UML/XML. J'ai dû résoudre les problèmes de parsing, d'injection de contexte dans les LLM et de respect des standards architecturaux. Je suis donc immédiatement opérationnel sur ce volet technique.

\textbf{2. Je maîtrise les architectures IA agentiques (Mission Catalogue \& Déploiement).}
En concevant un \textbf{système multi-agents agricole} (Google ADK, LangChain, RAG), j'ai expérimenté l'orchestration d'agents autonomes avec gestion de la mémoire contextuelle et résolution de conflits. Cette expérience me permet d'identifier les cas d'usage pertinents et les évolutions nécessaires pour intégrer l'IA dans les pratiques d'ingénierie système.

\textbf{3. J'ai une rigueur scientifique pour l'évaluation des solutions IA.}
Mon expérience d'Assistant de Recherche au \textbf{Mila (Institut Québécois d'IA)} sur l'interprétabilité des modèles m'a formé à l'analyse critique. J'y ai implémenté des pipelines de tests statistiques sous PyTorch et reproduit des expériences SOTA. Cette rigueur sera un atout pour évaluer les démonstrateurs existants et identifier les solutions différenciantes.

En parallèle, mon expérience en développement backend (Laravel, Docker, CI/CD) garantit que je produirai des outils intégrables dans vos environnements de production, pas uniquement du code de recherche.

Disponible dès avril 2026, je serais ravi de vous présenter mes projets lors d'un entretien.

Cordialement,

\vspace{0.5cm}

\textbf{Loïc Aron MBASSI EWOLO}

\end{document}
